\section{Introduction to the Link Layer (6.1)}
%=============================================================================

\subsection{Terminology and Context}

\begin{keybox}[Link Layer -- Basic Terminology]
\begin{itemize}[noitemsep]
    \item \textbf{Nodes:} Hosts and routers (any device running a link-layer protocol).
    \item \textbf{Links:} Communication channels that connect adjacent nodes along a communication path. Links can be wired (e.g., Ethernet cable), wireless (e.g., WiFi), or LANs.
    \item \textbf{Frame:} The link-layer PDU (protocol data unit). The link layer encapsulates a network-layer datagram into a frame.
\end{itemize}

\medskip
The link layer is responsible for transferring a datagram from one node to a \textbf{physically adjacent} node over a single link.
\end{keybox}

\begin{warnbox}[Important Distinction]
Different link protocols may be used on different links along the same path. For example, a datagram might travel over Ethernet on the first link, then over a point-to-point link, then over WiFi -- each hop uses whatever link-layer protocol is appropriate for that link.
\end{warnbox}

\subsection{Link Layer Services}

\begin{keybox}[Services Provided by the Link Layer]
\begin{enumerate}[noitemsep]
    \item \textbf{Framing:} Encapsulate datagram into a frame, adding header and trailer.
    \item \textbf{Link access (MAC protocol):} Controls how frames are transmitted onto the link. Especially important for shared (broadcast) links, where multiple nodes share the same medium.
    \item \textbf{Reliable delivery:} Guarantees to move each datagram across the link without error (not all link-layer protocols offer this; common on high-error-rate links like wireless).
    \item \textbf{Flow control:} Pacing between adjacent sending and receiving nodes.
    \item \textbf{Error detection:} Errors caused by signal attenuation and electromagnetic noise. The receiver detects the presence of errors.
    \item \textbf{Error correction:} The receiver identifies \emph{and corrects} bit errors without retransmission.
    \item \textbf{Half-duplex and full-duplex:} Half-duplex means nodes at both ends can transmit, but not at the same time.
\end{enumerate}
\end{keybox}

\begin{imagebox}[Link Layer Context]
\textbf{Textbook:} Kurose \& Ross 8e, Figure 6.1\\
\textbf{Google:} ``link layer protocol stack context nodes links diagram''\\
\textit{Shows where the link layer sits in the protocol stack and how it connects adjacent nodes over individual links.}
\end{imagebox}

\subsection{Where is the Link Layer Implemented?}

\begin{keybox}[Network Interface Card (NIC)]
The link layer is implemented in the \textbf{Network Interface Card (NIC)}, also known as the \textbf{network adapter}:
\begin{itemize}[noitemsep]
    \item Implements link-layer and physical-layer functionality
    \item Attaches to the host's system bus (PCIe, USB, etc.)
    \item Is a combination of hardware, software, and firmware
    \item The ``adaptor'' communicates with the rest of the host system
\end{itemize}

\medskip
\textbf{Sending side:} Encapsulates datagram in a frame, adds error-checking bits, manages flow control, etc.\\
\textbf{Receiving side:} Looks for errors, extracts datagram, passes it up to the receiving node's network layer.
\end{keybox}

%=============================================================================
\section{Error Detection and Correction (6.2)}
%=============================================================================

\subsection{Overview}

\begin{keybox}[Error Detection and Correction -- EDC]
Error detection is \textbf{not 100\% reliable}: the protocol may miss some errors (i.e., undetected bit errors). However, we can make the probability of undetected errors extremely small by using stronger (larger) EDC fields.

The general approach:
\begin{itemize}[noitemsep]
    \item Sender computes EDC bits based on the data $D$
    \item Both $D$ and EDC are sent
    \item Receiver recomputes EDC over received $D'$ and compares with received EDC
\end{itemize}
\end{keybox}

\subsection{Parity Checking}

\begin{keybox}[Single-Bit Parity]
Add a single parity bit so that the total number of 1's is even (even parity) or odd (odd parity).
\begin{itemize}[noitemsep]
    \item Can detect \textbf{single-bit} errors
    \item Cannot detect errors where an even number of bits are flipped
\end{itemize}
\end{keybox}

\begin{keybox}[Two-Dimensional (2D) Bit Parity]
Arrange the data bits in a matrix (rows and columns). Compute parity for each row and each column:
\begin{itemize}[noitemsep]
    \item Can \textbf{detect and correct} single-bit errors (identify the exact row and column of the flipped bit)
    \item Can \textbf{detect} (but not correct) any combination of two-bit errors
\end{itemize}
\end{keybox}

\begin{exambox}[Exam-Type: 2D Even Parity -- V23 Q1.4.2, V24 Q1.4.2]
\textbf{Question:} Which alternatives show correct implementation(s) of a two-dimensional even parity scheme?

\medskip
You are given 8 matrices (a--h), each with 4 rows of 4 data bits plus a parity column and a parity row. To verify:
\begin{enumerate}[noitemsep]
    \item Count the 1's in each \textbf{row} (including the parity bit). Must be even.
    \item Count the 1's in each \textbf{column} (including the parity bit). Must be even.
    \item If all rows AND all columns have even counts $\Rightarrow$ correct.
\end{enumerate}

\textbf{Tip:} Carefully check \emph{every} row and column. A single mismatch means that alternative is wrong.
\end{exambox}

\subsection{Internet Checksum}

\begin{keybox}[Internet Checksum (Review)]
Used in the transport layer (TCP/UDP) and IP header:
\begin{itemize}[noitemsep]
    \item Treat data as a sequence of 16-bit integers
    \item Compute the 1's complement sum of all 16-bit words
    \item The checksum is the 1's complement of that sum
    \item Receiver adds all words including checksum; result should be \texttt{1111111111111111}
\end{itemize}

\textbf{Weakness:} Relatively weak -- can miss certain patterns of errors. Used at the transport/network layer but \emph{not} typically at the link layer, which uses CRC instead.
\end{keybox}

\subsection{Cyclic Redundancy Check (CRC)}

\begin{keybox}[CRC -- Key Concepts]
CRC is a more powerful error-detection technique used at the link layer:
\begin{itemize}[noitemsep]
    \item Based on \textbf{modulo-2 arithmetic} (XOR, no carries/borrows)
    \item Uses a \textbf{generator polynomial} $G$ of degree $r$ (i.e., $r+1$ bits long)
    \item The sender appends $r$ CRC bits (remainder $R$) to the data $D$
    \item The transmitted bits are: $D \cdot 2^r \oplus R$
    \item The receiver divides the received bits by $G$; if the remainder is 0, no error is detected
\end{itemize}
\end{keybox}

\begin{formulabox}[CRC Computation]
Given data bits $D$ and generator $G$ (of $r+1$ bits):
\begin{enumerate}[noitemsep]
    \item Append $r$ zero bits to $D$: $D \cdot 2^r$
    \item Divide $D \cdot 2^r$ by $G$ using modulo-2 division (XOR)
    \item The remainder $R$ (exactly $r$ bits) is the CRC
    \item Transmit: $D \cdot 2^r \oplus R$ (which is exactly divisible by $G$)
\end{enumerate}

\medskip
\textbf{Key properties:}
\begin{itemize}[noitemsep]
    \item Can detect \textbf{all burst errors of fewer than $r+1$ bits}
    \item Can detect \textbf{all odd numbers of bit errors} (if $G$ contains $(x+1)$ as a factor)
    \item Burst errors of length $> r+1$ are detected with probability $1 - (1/2)^r$
\end{itemize}
\end{formulabox}

\begin{imagebox}[Error Detection and CRC]
\textbf{Textbook:} Kurose \& Ross 8e, Figure 6.4 and Figure 6.6\\
\textbf{Google:} ``CRC cyclic redundancy check calculation example diagram''\\
\textit{Illustrates the error-detection scheme and step-by-step CRC modulo-2 division.}
\end{imagebox}

\begin{warnbox}[CRC vs.\ Checksum]
CRC is much stronger than the Internet checksum. CRC is used at the link layer (e.g., Ethernet CRC-32), while the Internet checksum is used at higher layers. CRC can detect all burst errors up to length $r$, whereas the checksum can miss many error patterns.
\end{warnbox}

\begin{exambox}[Exam-Type: CRC Matching -- S2025 Q7]
\textbf{Question:} Given data string 101110 and a CRC generator, match the following:
\begin{itemize}[noitemsep]
    \item CRC value $\rightarrow$ 001
    \item CRC length $\rightarrow$ 3
    \item CRC generator $\rightarrow$ 1001
    \item Total bits sent $\rightarrow$ 9
    \item CRC generator polynomial $\rightarrow$ $x^3 + 1$
\end{itemize}

\textbf{How to solve:} The generator 1001 has 4 bits, so $r = 3$. Append 3 zeros to data: 101110\textbf{000}. Perform modulo-2 division by 1001 to get remainder 001. Total bits = 6 (data) + 3 (CRC) = 9. The polynomial $1001 = 1\cdot x^3 + 0\cdot x^2 + 0\cdot x^1 + 1\cdot x^0 = x^3 + 1$.
\end{exambox}

%=============================================================================
\section{Multiple Access Protocols (6.3)}
%=============================================================================

\subsection{The Multiple Access Problem}

\begin{keybox}[The Shared Channel Problem]
On a \textbf{broadcast link} (shared medium), when two or more nodes transmit simultaneously, a \textbf{collision} occurs -- the signals interfere and the transmitted frames are corrupted.

\textbf{Goal of MAC (Medium Access Control) protocols:} Determine how nodes share the channel, i.e., determine when a node can transmit.

\medskip
\textbf{Ideal properties for a MAC protocol} (channel rate $R$):
\begin{enumerate}[noitemsep]
    \item When one node wants to transmit, it can send at rate $R$
    \item When $M$ nodes want to transmit, each gets average rate $R/M$
    \item Fully decentralized (no master node, no clock synchronization)
    \item Simple
\end{enumerate}
\end{keybox}

\subsection{Taxonomy of MAC Protocols}

\begin{keybox}[Three Categories of MAC Protocols]
\begin{enumerate}
    \item \textbf{Channel Partitioning:} Divide channel into smaller ``pieces'' (time slots, frequency bands, codes). Allocate a piece to each node for exclusive use.
    \begin{itemize}[noitemsep]
        \item TDMA, FDMA, CDMA
    \end{itemize}

    \item \textbf{Random Access:} Channel not divided; allow collisions, then ``recover'' from them.
    \begin{itemize}[noitemsep]
        \item ALOHA, Slotted ALOHA, CSMA, CSMA/CD, CSMA/CA
    \end{itemize}

    \item \textbf{Taking Turns:} Nodes take turns, but nodes with more to send can take longer turns.
    \begin{itemize}[noitemsep]
        \item Polling, Token passing
    \end{itemize}
\end{enumerate}
\end{keybox}

\subsection{Channel Partitioning Protocols}

\begin{keybox}[TDMA -- Time Division Multiple Access]
\begin{itemize}[noitemsep]
    \item Access to the channel is divided into ``rounds'' of fixed-length time slots
    \item Each station gets a fixed-length slot in every round
    \item Unused slots go idle (wasted capacity)
    \item \textbf{Pros:} No collisions, fair
    \item \textbf{Cons:} Wasted capacity when nodes have nothing to send; each node limited to $R/N$ even if it's the only one transmitting
\end{itemize}
\end{keybox}

\begin{keybox}[FDMA -- Frequency Division Multiple Access]
\begin{itemize}[noitemsep]
    \item Channel spectrum divided into frequency bands
    \item Each station gets a fixed frequency band
    \item Unused bands are idle (wasted capacity)
    \item Same pros/cons as TDMA
\end{itemize}
\end{keybox}

\subsection{Random Access Protocols}

\subsubsection{Slotted ALOHA}

\begin{keybox}[Slotted ALOHA]
\textbf{Assumptions:}
\begin{itemize}[noitemsep]
    \item All frames are exactly the same size
    \item Time is divided into equal-size slots (= time to transmit 1 frame)
    \item Nodes start to transmit only at the beginning of a slot
    \item Nodes are synchronized (they all know when slots begin)
    \item If 2 or more nodes transmit in the same slot, all nodes detect the collision
\end{itemize}

\textbf{Operation:}
\begin{itemize}[noitemsep]
    \item When a node has a frame to send, it transmits in the next slot
    \item If no collision: success!
    \item If collision: retransmit in each subsequent slot with probability $p$, until success
\end{itemize}
\end{keybox}

\begin{formulabox}[Slotted ALOHA Efficiency]
Maximum efficiency = $1/e \approx 0.37$ (37\%)

This means: even under best conditions, the channel is productive only 37\% of the time. The remaining 63\% is wasted on collisions and empty slots.

\medskip
Derivation: With $N$ active nodes, each transmitting with probability $p$, the probability of a successful slot is $N \cdot p \cdot (1-p)^{N-1}$. Optimizing over $p$ and letting $N \to \infty$ gives $\lim_{N\to\infty} N \cdot \frac{1}{N} \cdot (1 - \frac{1}{N})^{N-1} = 1/e$.
\end{formulabox}

\subsubsection{Pure (Unslotted) ALOHA}

\begin{keybox}[Pure ALOHA]
\begin{itemize}[noitemsep]
    \item No synchronization, no slots
    \item A node transmits immediately when it has a frame
    \item Collision risk: if any other node transmits during the \emph{entire} vulnerable period (2 frame times), there's a collision
    \item \textbf{Maximum efficiency = $1/(2e) \approx 0.18$ (18\%)} -- half that of slotted ALOHA
\end{itemize}
\end{keybox}

\begin{exambox}[Exam-Type: ALOHA Properties -- V24 Q1.4.5]
\textbf{Question:} Which statements about pure ALOHA and slotted ALOHA are true?

\textbf{Key facts:}
\begin{itemize}[noitemsep]
    \item Pure ALOHA has lower efficiency than slotted ALOHA \checkmark
    \item Max efficiency of slotted ALOHA ($1/e$) is \textbf{two times} that of pure ALOHA ($1/2e$) \checkmark
    \item Slotted ALOHA reduces collisions compared to pure ALOHA \checkmark
    \item In slotted ALOHA, stations must wait for the beginning of the next time slot $\Rightarrow$ they \textbf{cannot} send at any time without synchronization
    \item In pure ALOHA, stations can transmit at any time (no waiting for slots)
\end{itemize}
\end{exambox}

\subsubsection{CSMA -- Carrier Sense Multiple Access}

\begin{keybox}[CSMA -- ``Listen Before Transmit'']
\textbf{CSMA} adds \textbf{carrier sensing}: if the channel is sensed busy, defer transmission.

\medskip
\textbf{Why do collisions still occur?} Because of \textbf{propagation delay}: a node may not hear another node's transmission that has already started (the signal hasn't arrived yet). The longer the propagation delay, the higher the chance of collision.
\end{keybox}

\subsubsection{CSMA/CD -- CSMA with Collision Detection}

\begin{keybox}[CSMA/CD (used in Ethernet)]
CSMA/CD adds \textbf{collision detection}: if a collision is detected, the transmission is \textbf{aborted immediately}, reducing channel waste.

\textbf{Algorithm:}
\begin{enumerate}[noitemsep]
    \item NIC receives datagram from network layer, creates frame
    \item If NIC senses channel idle, start transmitting. If busy, wait until idle, then transmit.
    \item If entire frame transmitted without detecting collision $\Rightarrow$ done!
    \item If collision detected, abort and send a \textbf{jam signal}
    \item After aborting, enter \textbf{binary (exponential) backoff}
\end{enumerate}
\end{keybox}

\begin{formulabox}[Binary Exponential Backoff]
After the $m$th collision (for the same frame):
\begin{itemize}[noitemsep]
    \item Choose $K$ randomly from $\{0, 1, 2, \ldots, 2^m - 1\}$
    \item Wait $K \times 512$ bit times (for Ethernet)
    \item After 10 collisions, freeze the range at $\{0, 1, \ldots, 1023\}$
    \item After 16 collisions, give up and report an error
\end{itemize}

The idea: adapt the retransmission delay to the estimated number of colliding stations. More collisions $\Rightarrow$ longer expected backoff.
\end{formulabox}

\begin{imagebox}[CSMA/CD Operation]
\textbf{Textbook:} Kurose \& Ross 8e, Figure 6.12\\
\textbf{Google:} ``CSMA CD collision detection Ethernet algorithm diagram''\\
\textit{Shows the CSMA/CD algorithm: sense, transmit, detect collision, jam, and exponential backoff.}
\end{imagebox}

\begin{formulabox}[CSMA/CD Efficiency]
\[
\text{Efficiency} = \frac{1}{1 + 5 \cdot t_{\text{prop}} / t_{\text{trans}}}
\]
where:
\begin{itemize}[noitemsep]
    \item $t_{\text{prop}}$ = maximum propagation delay between any 2 nodes
    \item $t_{\text{trans}}$ = time to transmit a maximum-size frame
\end{itemize}

\medskip
Key insight: Efficiency $\to 1$ as $t_{\text{prop}} \to 0$ or $t_{\text{trans}} \to \infty$.\\
Better than ALOHA, and is decentralized, simple, and cheap.
\end{formulabox}

\subsection{Taking Turns Protocols}

\begin{keybox}[Polling and Token Passing]
\textbf{Polling:}
\begin{itemize}[noitemsep]
    \item A master node ``invites'' slave nodes to transmit in turn
    \item \textbf{Concerns:} Polling overhead, latency, single point of failure (master)
    \item Example: Bluetooth
\end{itemize}

\textbf{Token Passing:}
\begin{itemize}[noitemsep]
    \item A special ``token'' frame is passed from node to node sequentially
    \item A node can transmit only when it holds the token
    \item \textbf{Concerns:} Token overhead, latency, single point of failure (token loss)
    \item Example: FDDI (Fiber Distributed Data Interface)
\end{itemize}
\end{keybox}

\begin{exambox}[Exam-Type: Classify MAC Protocols -- S2025 Q6]
\textbf{Question:} Categorize these protocols into Channel Partitioning, Random Access, or Taking Turns:

\begin{center}
\begin{tabular}{ll}
\toprule
\textbf{Protocol} & \textbf{Category} \\
\midrule
Bluetooth & Taking Turns (polling) \\
Ethernet CSMA/CD & Random Access \\
CDMA & Channel Partitioning \\
Slotted ALOHA & Random Access \\
FDDI & Taking Turns (token passing) \\
FDMA & Channel Partitioning \\
\bottomrule
\end{tabular}
\end{center}
\end{exambox}

\subsection{Cable Access Network: DOCSIS}

\begin{keybox}[DOCSIS -- Data Over Cable Service Interface Specification]
Cable internet uses a shared medium (coaxial cable + fiber):
\begin{itemize}[noitemsep]
    \item \textbf{Downstream:} FDM (Frequency Division Multiplexing) -- multiple channels, broadcast to all homes
    \item \textbf{Upstream:} FDM + TDM + Random Access
    \item Upstream channels are divided into time slots; some slots are assigned (for data), some are for contention (random access using mini-slots)
    \item MAP messages from CMTS (Cable Modem Termination System) assign upstream slots
    \item Random access slots use a form of ALOHA with binary backoff for requesting bandwidth
\end{itemize}
\end{keybox}

%=============================================================================
\section{LANs -- Addressing and ARP (6.4)}
%=============================================================================

\subsection{MAC Addresses}

\begin{keybox}[MAC (Physical/LAN) Addresses]
\begin{itemize}[noitemsep]
    \item \textbf{48-bit} address (6 bytes), written in hexadecimal (e.g., \texttt{1A-2F-BB-76-09-AD})
    \item \textbf{Burned into} the NIC's ROM (or set by software)
    \item Administered by IEEE: manufacturers buy blocks of addresses
    \item \textbf{Flat address space:} portable -- you can move a NIC from one LAN to another without changing its MAC address
    \item \textbf{Broadcast address:} \texttt{FF-FF-FF-FF-FF-FF} -- delivered to all adapters on the LAN
\end{itemize}
\end{keybox}

\begin{warnbox}[MAC vs.\ IP Addresses]
\begin{center}
\begin{tabular}{lll}
\toprule
& \textbf{MAC Address} & \textbf{IP Address} \\
\midrule
Length & 48 bits (6 bytes) & 32 bits (IPv4) \\
Structure & Flat (no hierarchy) & Hierarchical (network + host) \\
Scope & Link-layer, local & Network-layer, global \\
Portability & Portable (burned in NIC) & Depends on subnet \\
Analogy & SSN (identity number) & Postal address (location) \\
\bottomrule
\end{tabular}
\end{center}

\textbf{Why both?} IP addresses are needed for network-layer routing (hierarchical, location-dependent). MAC addresses are needed for link-layer frame delivery on a single LAN segment.
\end{warnbox}

\subsection{ARP -- Address Resolution Protocol}

\begin{keybox}[ARP -- How It Works]
ARP maps an IP address to a MAC address \emph{on the same subnet}:

\begin{enumerate}[noitemsep]
    \item Node A wants to send a frame to Node B (knows B's IP, but not B's MAC)
    \item A broadcasts an \textbf{ARP query} (dest MAC = \texttt{FF-FF-FF-FF-FF-FF}) containing B's IP address
    \item \textbf{All nodes} on the LAN receive the query
    \item Node B recognizes its IP and replies with a \textbf{unicast ARP response} containing its MAC address
    \item A caches the IP$\to$MAC mapping in its \textbf{ARP table} (with TTL, typically 20 minutes)
\end{enumerate}

ARP is ``plug-and-play'': tables are built automatically, no manual configuration needed.
\end{keybox}

\begin{exambox}[Exam-Type: ARP -- V23 Q1.4.1]
\textbf{Question:} It is \_\_\_\_'s job to translate between IP addresses and MAC (medium access control) addresses.

\textbf{Answer:} \textbf{ARP (Address Resolution Protocol)}.

\medskip
\textbf{Common confusion:}
\begin{itemize}[noitemsep]
    \item \textbf{DNS} translates \emph{hostnames} to \emph{IP addresses} (application layer)
    \item \textbf{ARP} translates \emph{IP addresses} to \emph{MAC addresses} (link layer)
\end{itemize}
\end{exambox}

\begin{imagebox}[ARP Operation]
\textbf{Textbook:} Kurose \& Ross 8e, Figure 6.19\\
\textbf{Google:} ``ARP address resolution protocol how it works diagram''\\
\textit{Shows the ARP table and the broadcast query / unicast reply process for resolving IP to MAC.}
\end{imagebox}

\subsection{Routing to Another Subnet}

\begin{keybox}[Sending a Datagram to Another Subnet]
When host A (on subnet 1) sends to host B (on subnet 2), going through router R:

\textbf{Key principle:} \textbf{IP addresses stay the same} end-to-end; \textbf{MAC addresses change at each hop}.

\begin{enumerate}[noitemsep]
    \item A creates IP datagram: src IP = A, dest IP = B
    \item A looks up next-hop (router R) using its routing table
    \item A uses ARP to find R's MAC address on subnet 1
    \item A creates link-layer frame: src MAC = A, dest MAC = R (on subnet 1 interface)
    \item Router R receives frame, extracts datagram, looks up next hop for B
    \item R uses ARP to find B's MAC address on subnet 2
    \item R creates new frame: src MAC = R (on subnet 2 interface), dest MAC = B
    \item B receives the frame
\end{enumerate}
\end{keybox}

\begin{warnbox}[Critical Exam Point: MAC Changes -- IP Stays]
At every hop:
\begin{itemize}[noitemsep]
    \item The \textbf{source and destination IP addresses} remain \textbf{unchanged} (A $\to$ B)
    \item The \textbf{source and destination MAC addresses} are \textbf{rewritten} to reflect the current link
\end{itemize}
This is a very commonly tested concept!
\end{warnbox}

%=============================================================================
\section{Ethernet (6.4)}
%=============================================================================

\subsection{Ethernet Overview}

\begin{keybox}[Ethernet -- The Dominant Wired LAN Technology]
Ethernet is the most widely used wired LAN technology:
\begin{itemize}[noitemsep]
    \item First widely used LAN technology
    \item Kept up with speed demands: 10 Mbps $\to$ 10 Gbps $\to$ 400 Gbps
    \item Simple, cheap, and standardized (IEEE 802.3)
\end{itemize}
\end{keybox}

\subsection{Ethernet Frame Structure}

\begin{keybox}[Ethernet Frame Format]
\begin{center}
\begin{tabular}{|c|c|c|c|c|c|}
\hline
\textbf{Preamble} & \textbf{Dest MAC} & \textbf{Src MAC} & \textbf{Type} & \textbf{Data (Payload)} & \textbf{CRC} \\
\hline
8 bytes & 6 bytes & 6 bytes & 2 bytes & 46--1500 bytes & 4 bytes \\
\hline
\end{tabular}
\end{center}

\begin{itemize}[noitemsep]
    \item \textbf{Preamble} (8 bytes): 7 bytes of \texttt{10101010} followed by 1 byte of \texttt{10101011} (SFD -- Start Frame Delimiter). Used for clock synchronization.
    \item \textbf{Destination MAC} (6 bytes): MAC address of the receiving adapter. If broadcast (\texttt{FF-FF-FF-FF-FF-FF}), all adapters process the frame.
    \item \textbf{Source MAC} (6 bytes): MAC address of the sending adapter.
    \item \textbf{Type} (2 bytes): Indicates the higher-layer protocol (e.g., \texttt{0x0800} = IPv4, \texttt{0x0806} = ARP). Acts as a demultiplexing key.
    \item \textbf{Data/Payload} (46--1500 bytes): The network-layer datagram. Minimum 46 bytes (padded if necessary). Maximum 1500 bytes = \textbf{MTU} (Maximum Transmission Unit).
    \item \textbf{CRC} (4 bytes): Cyclic Redundancy Check for error detection. If CRC check fails, the frame is \textbf{dropped} (no NAK/retransmission at link layer).
\end{itemize}
\end{keybox}

\begin{imagebox}[Ethernet Frame Structure]
\textbf{Textbook:} Kurose \& Ross 8e, Figure 6.20\\
\textbf{Google:} ``Ethernet frame format structure fields diagram''\\
\textit{Shows the Ethernet frame fields: preamble, destination/source MAC, type, payload, and CRC.}
\end{imagebox}

\subsection{Ethernet Properties}

\begin{keybox}[Ethernet -- Key Properties]
\begin{itemize}[noitemsep]
    \item \textbf{Connectionless:} No handshaking between sending and receiving NICs.
    \item \textbf{Unreliable:} Receiving NIC does \textbf{not} send ACKs or NAKs. Dropped frames are only recovered if the sender uses a reliable higher-layer protocol (e.g., TCP).
    \item \textbf{MAC protocol:} Uses unslotted CSMA/CD with binary exponential backoff (in half-duplex mode).
\end{itemize}
\end{keybox}

\subsection{Ethernet Topologies}

\begin{keybox}[Bus vs.\ Star Topology]
\textbf{Bus topology (legacy):}
\begin{itemize}[noitemsep]
    \item All nodes share a single coaxial cable (broadcast medium)
    \item Collisions are common; uses CSMA/CD
    \item Popular in the 1990s
\end{itemize}

\textbf{Star topology (modern):}
\begin{itemize}[noitemsep]
    \item Each node has a dedicated connection to a central \textbf{switch}
    \item The switch operates in full-duplex mode -- \textbf{no collisions!}
    \item Each link is its own collision domain
    \item Dominant topology today
\end{itemize}
\end{keybox}

\subsection{802.3 Ethernet Standards}

\begin{keybox}[Common Ethernet Standards]
\begin{center}
\begin{tabular}{lll}
\toprule
\textbf{Standard} & \textbf{Speed} & \textbf{Medium} \\
\midrule
100BASE-TX & 100 Mbps & Copper (Cat5) \\
1000BASE-T & 1 Gbps & Copper (Cat5e/6) \\
10GBASE-T & 10 Gbps & Copper (Cat6a) \\
100GBASE-SR & 100 Gbps & Fiber \\
400GBASE & 400 Gbps & Fiber \\
\bottomrule
\end{tabular}
\end{center}

All share the same frame format and use CSMA/CD (at lower speeds) or full-duplex switched (at higher speeds).
\end{keybox}

%=============================================================================
\section{Link-Layer Switches (6.4)}
%=============================================================================

\subsection{Switch Basics}

\begin{keybox}[Ethernet Switch -- Key Properties]
An Ethernet switch is a \textbf{link-layer device}:
\begin{itemize}[noitemsep]
    \item \textbf{Store-and-forward:} Examines incoming frame's destination MAC, selectively forwards to one or more outgoing links
    \item \textbf{Transparent:} Hosts are \textbf{unaware} of the switch's presence
    \item \textbf{Plug-and-play / self-learning:} Switches do \textbf{not} need to be manually configured
    \item Uses a \textbf{switch table}: entries of (MAC address, interface, TTL)
\end{itemize}
\end{keybox}

\subsection{Switch Self-Learning}

\begin{keybox}[Self-Learning Algorithm]
The switch \textbf{learns} which hosts are reachable through which interfaces:
\begin{enumerate}[noitemsep]
    \item When a frame arrives on interface $i$, the switch records the \textbf{source MAC address} in its switch table, along with the interface $i$ and a timestamp
    \item If an entry is not refreshed within the \textbf{TTL} (aging time), it is deleted
\end{enumerate}

This is completely automatic -- no configuration needed!
\end{keybox}

\begin{imagebox}[Switch Self-Learning]
\textbf{Textbook:} Kurose \& Ross 8e, Figure 6.22\\
\textbf{Google:} ``Ethernet switch self-learning forwarding table diagram''\\
\textit{Shows an institutional network with interconnected switches and how switch tables are populated.}
\end{imagebox}

\subsection{Switch Forwarding}

\begin{keybox}[Frame Forwarding/Filtering]
When a frame arrives at a switch on interface $i$, with destination MAC address \texttt{dest}:

\begin{enumerate}[noitemsep]
    \item Look up \texttt{dest} in the switch table
    \item \textbf{If entry found:}
    \begin{itemize}[noitemsep]
        \item If \texttt{dest} is on interface $i$ (same segment) $\Rightarrow$ \textbf{drop/filter} (no need to forward)
        \item If \texttt{dest} is on interface $j \neq i$ $\Rightarrow$ \textbf{forward} the frame to interface $j$
    \end{itemize}
    \item \textbf{If entry not found:} \textbf{flood} -- forward the frame on \textbf{all} interfaces except $i$
\end{enumerate}
\end{keybox}

\subsection{Switches vs.\ Routers}

\begin{keybox}[Switches vs.\ Routers]
\begin{center}
\begin{tabular}{lll}
\toprule
& \textbf{Switch (Layer 2)} & \textbf{Router (Layer 3)} \\
\midrule
Layer & Data link layer & Network layer \\
Addresses & MAC addresses & IP addresses \\
Forwarding & Switch table (self-learned) & Routing table (algorithms) \\
Broadcast & Floods unknown frames & Does not forward broadcasts \\
Configuration & Plug-and-play & Needs IP configuration \\
Scalability & Limited (flat addressing) & Hierarchical (IP subnets) \\
\bottomrule
\end{tabular}
\end{center}
\end{keybox}

\begin{imagebox}[Switches vs.\ Routers]
\textbf{Textbook:} Kurose \& Ross 8e, Figure 6.26\\
\textbf{Google:} ``switch vs router comparison network layers diagram''\\
\textit{Shows how packets flow through switches (layer 2) and routers (layer 3) in a network.}
\end{imagebox}

\begin{exambox}[Exam-Type: Switches vs.\ Routers -- V24 Q1.4.1]
\textbf{Question:} What is/are \textbf{false} when comparing between switches and routers?

\textbf{Key facts:}
\begin{itemize}[noitemsep]
    \item Both routers and switches are connecting devices in networking \checkmark
    \item Routers operate at the \textbf{Data link layer} and switches at \textbf{Network layer} $\Rightarrow$ \textbf{FALSE!} (It's the opposite)
    \item Switches operate at the Data link layer, routers at the Network layer \checkmark
    \item Switches connect devices within a network; routers connect different networks \checkmark
    \item Routers rely on IP addresses; switches rely on MAC addresses \checkmark
\end{itemize}
\end{exambox}

\begin{imagebox}[VLANs]
\textbf{Textbook:} Kurose \& Ross 8e, Figure 6.25\\
\textbf{Google:} ``VLAN virtual LAN trunk port switch diagram''\\
\textit{Shows how VLANs partition a single physical switch into multiple virtual LANs with trunk links.}
\end{imagebox}

%=============================================================================
\section{A Day in the Life of a Web Request (6.7)}
%=============================================================================

\begin{keybox}[Synthesis: All Layers Working Together]
When a laptop connects to a network and requests a web page, the following protocols interact across \textbf{all five layers}:

\textbf{Step 1: Connecting to the network (DHCP)}
\begin{itemize}[noitemsep]
    \item Laptop needs an IP address, subnet mask, default gateway, and DNS server
    \item DHCP Discover $\to$ DHCP Offer $\to$ DHCP Request $\to$ DHCP ACK
    \item DHCP is encapsulated in UDP, then IP, then Ethernet (broadcast)
\end{itemize}

\textbf{Step 2: Getting the router's MAC address (ARP)}
\begin{itemize}[noitemsep]
    \item Laptop knows the gateway IP (from DHCP) but needs its MAC address
    \item ARP request (broadcast) $\to$ ARP reply (unicast) gives the router's MAC
\end{itemize}

\textbf{Step 3: Resolving the web server's IP (DNS)}
\begin{itemize}[noitemsep]
    \item Laptop creates DNS query for the web server hostname
    \item DNS query sent via UDP to DNS server (IP from DHCP)
    \item DNS response provides the web server's IP address
\end{itemize}

\textbf{Step 4: Establishing a connection (TCP)}
\begin{itemize}[noitemsep]
    \item TCP SYN to web server $\to$ SYN-ACK $\to$ ACK (three-way handshake)
    \item IP datagram routed through network; at each hop, link-layer frames are created
\end{itemize}

\textbf{Step 5: Requesting the web page (HTTP)}
\begin{itemize}[noitemsep]
    \item HTTP GET request sent inside TCP segment
    \item Web server responds with HTTP response containing the web page
\end{itemize}
\end{keybox}

\begin{exambox}[Exam-Type: Protocols at Each Step -- S2025 Part II Q3]
\textbf{Question:} A researcher plugs his laptop into the Ethernet port. He sends an email using Outlook. Elaborate the protocols used at each step.

\textbf{Answer outline:}
\begin{enumerate}[noitemsep]
    \item \textbf{DHCP} (over UDP/IP/Ethernet broadcast): Obtain IP address, gateway, DNS server
    \item \textbf{ARP}: Resolve gateway router's MAC address
    \item \textbf{DNS} (over UDP): Resolve mail server hostname to IP
    \item \textbf{TCP}: Establish connection to mail server
    \item \textbf{SMTP}: Send the email (application layer)
    \item On the receiving end: \textbf{IMAP/POP3/HTTP} to retrieve the email
\end{enumerate}

At every hop along the path, the \textbf{link layer} (Ethernet or other) handles frame creation and MAC addressing, while \textbf{IP} handles routing.
\end{exambox}

%=============================================================================
\section{Summary of Key Formulas and Facts}
%=============================================================================

\begin{formulabox}[Key Formulas -- Chapter 6]
\begin{center}
\begin{longtable}{p{5.5cm}p{8.5cm}}
\toprule
\textbf{Concept} & \textbf{Formula / Value} \\
\midrule
Slotted ALOHA max efficiency & $1/e \approx 37\%$ \\
\addlinespace
Pure ALOHA max efficiency & $1/(2e) \approx 18\%$ \\
\addlinespace
CSMA/CD efficiency & $\dfrac{1}{1 + 5 \cdot t_{\text{prop}} / t_{\text{trans}}}$ \\
\addlinespace
Binary exponential backoff & After $m$th collision: $K \in \{0, 1, \ldots, 2^m - 1\}$, wait $K \times 512$ bit times \\
\addlinespace
CRC remainder bits & $r = \deg(G) = $ length of $G$ minus 1 \\
\addlinespace
CRC burst detection & Detects all bursts $< r+1$ bits \\
\addlinespace
MAC address length & 48 bits (6 bytes) \\
\addlinespace
ARP table TTL & Typically 20 minutes \\
\addlinespace
Ethernet MTU & 1500 bytes \\
\addlinespace
Ethernet min payload & 46 bytes \\
\bottomrule
\end{longtable}
\end{center}
\end{formulabox}

%=============================================================================
\section{Additional Exam Questions from Past Papers}
%=============================================================================

\begin{exambox}[S2025 -- CSMA Analysis -- Part II Q7]
\textbf{Question:} Consider 6 packets arriving at times $t=0.1, 1.4, 1.8, 3.2, 3.3, 4.1$ at different wireless nodes. Each transmission takes exactly 1 time unit. Propagation delay = 0.2 time units. For CSMA (without collision detection), determine which packets are successfully transmitted.

\medskip
\textbf{Approach:}
\begin{enumerate}[noitemsep]
    \item Packet 1 starts at $t=0.1$, finishes at $t=1.1$. No other node has transmitted yet $\Rightarrow$ \textbf{success}.
    \item Packet 2 at $t=1.4$: channel is free (packet 1 ended at $t=1.1$). Check if signal from packet 2 reaches other nodes before they transmit.
    \item Packet 3 at $t=1.8$: does it sense packet 2? Packet 2 started at $t=1.4$, prop delay = 0.2, so signal reaches at $t=1.6$. At $t=1.8$, packet 3's node senses busy $\Rightarrow$ defers.
    \item Continue this analysis for each packet\ldots
\end{enumerate}

\textbf{Key points:}
\begin{itemize}[noitemsep]
    \item A node senses the channel \emph{before} transmitting
    \item Signal takes 0.2 time units to propagate between nodes
    \item If the channel is sensed busy, the node does \emph{not} transmit
    \item If collision occurs, both frames are lost (no CD, so full transmission happens)
\end{itemize}
\end{exambox}

\begin{exambox}[S2025 -- Link-Layer Frame Headers -- Q8]
\textbf{Question:} A link-layer frame travels from a host to a router. Match headers H$_1$, H$_2$, H$_3$ (outermost to innermost before the message M) to the correct layer.

\textbf{Answer:}
\begin{itemize}[noitemsep]
    \item H$_1$ (outermost) = \textbf{Link layer} (Ethernet header with MAC addresses)
    \item H$_2$ = \textbf{Network layer} (IP header)
    \item H$_3$ (closest to M) = \textbf{Transport layer} (TCP/UDP header)
\end{itemize}

Remember: encapsulation adds headers from top to bottom. At the link layer, the frame wraps everything, so the link-layer header is the outermost.
\end{exambox}

\begin{exambox}[V23/V24 -- Flow Control in Link Layer -- V23 Q2.2]
\textbf{Question:} Besides Transport Layer, which of Physical Layer, Link Layer, and Network Layer has flow control concepts?

\textbf{Answer:} The \textbf{Link Layer} has flow control (pacing between adjacent sending and receiving nodes). This is analogous to TCP flow control but operates at the link level.

\medskip
The Link Layer also has error detection/correction (like transport layer's checksum), and reliable delivery on some link protocols (like WiFi 802.11).
\end{exambox}

\begin{exambox}[S2025 -- Subnetting and Broadcast Address -- Q5 and Q12]
\textbf{Q5:} Given 208.27.1.0/22, create 12 subnets. How many usable hosts per subnet?

/22 = 1024 addresses total. To get 12 subnets, need $\lceil\log_2 12\rceil = 4$ subnet bits $\Rightarrow$ /26.
Each /26 subnet has $2^6 - 2 = 62$ usable hosts. \textbf{Answer: b) 62}

\medskip
\textbf{Q12:} What is the broadcast address of prefix 172.18.16.0/21?

/21 means 21 network bits, 11 host bits. $2^{11} = 2048$ addresses.
$172.18.16.0 + 2047 = 172.18.23.255$. \textbf{Answer: b) 172.18.23.255}
\end{exambox}
